% Volume III: Attention and Publicity
% Part V. Publicity Without Comprehension
% Chapter 26: Flooding, Noise, and Plausible Distraction
% Spine: proof-spines/ch026-spine.md

\chapter{Flooding, Noise, and Plausible Distraction}
\label{ch:ch026}

Censorship and flooding are dual strategies: censorship decreases accessible evidence, while flooding decreases the effective signal-to-attention ratio. If relevant signal has mass \(S\), distracting material has mass \(N\), and attention is bounded, effective discoverability may scale as
\[
D \propto \frac{S}{S+N}.
\]
\ToyModel\ The point is structural rather than psychological: when noise grows, it consumes the same interpretive budget that serious inquiry requires.

The chapter studies spam, document dumps, duplicated footage, performative transparency, adversarial disclosure, and archives that are technically accessible but practically unusable. Flooding preserves deniability because nothing need be hidden; what matters is simply prevented from remaining coherent long enough to bite.
